\section{Implementation Details} % show best code and explain
\label{sec:implementation}
Explain how the proposed approach is implemented in the project. Mention the programming language, major
modules, communication mechanism, configuration files, and how nodes are started and monitored.

\subsection{Node Agent}
\label{subsec:node_agent}
Describe the responsibilities of the node agent, including sending or responding to health checks, reporting
status, and interacting with the membership control service.

\subsubsection{Hybrid SWIMPHI}
\label{subsubsec:swimphi_code}

The SWIMPHI is implemented by the \texttt{FailureDetector} class, which combines periodic heartbeat probing inspired by
SWIM with adaptive suspicion estimation from the $\phi$ Accrual Failure Detector.

As shown in Listing~\ref{lst:scheduler}, the \texttt{FailureDetector} uses a \texttt{ScheduledExecutorService} to
schedule periodic monitoring tasks. The primary task performs heartbeat probing at configurable intervals, enabling
each node to continuously monitor the availability of its neighbors and detect potential failures in a timely manner.

\begin{lstlisting}[language=Java,
caption={Scheduling periodic failure-detection tasks},
label={lst:scheduler}]
scheduler.scheduleAtFixedRate(
    this::runOneProbeSafely,
    probeIntervalSeconds,
    probeIntervalSeconds,
    TimeUnit.SECONDS);
\end{lstlisting}

During each monitoring cycle, the node selects a target neighbor from the \texttt{NeighborDirectory} and sends a
heartbeat request through the \texttt{NodeClient}. The implementation uses Java's \texttt{CompletableFuture} framework
to perform non-blocking communication, allowing monitoring operations to continue without waiting for network
responses.

\begin{lstlisting}[language=Java,
caption={Heartbeat probing and ACK processing},
label={lst:heartbeat}]
Optional<NodeAddress> selectedTarget = neighborDirectory.nextTargetNode();

nodeClient.ping(targetNode).thenAccept(ackReceived -> {
    if (ackReceived)
        handleAckReceived(
          targetNode, "direct ping", pingSendTime);
    else handleDirectPingFailure(targetNode);
});
\end{lstlisting}

When an acknowledgment (ACK) is received, the monitored node is marked as \textit{ALIVE}. The detector records the
heartbeat arrival time and updates a sliding window of historical heartbeat intervals maintained for that neighbor.
These observations are subsequently used by the $\phi$ Accrual Failure Detector to model normal communication behavior.

If a direct heartbeat probe does not receive an acknowledgment, the detector invokes the indirect verification
procedure. As shown in the implementation, a set of helper nodes is selected from the local neighbor table using the
\texttt{NeighborDirectory} component.

\begin{lstlisting}[language=Java,
caption={Selection of helper nodes for indirect verification},
label={lst:helper}]
List<NodeAddress> helperNodes = neighborDirectory.selectHelperNodes(targetNode);
\end{lstlisting}

The selected helper nodes are intended to perform \texttt{PING\_REQ} operations on behalf of tfhe original sender. Each
helper node attempts to contact the target node and reports the result back to the requester. If at least one helper
node successfully receives an acknowledgment, the target node is considered reachable and its heartbeat statistics are
updated accordingly.

If no acknowledgment is received in any helper nodes, the detector evaluates the target node using the $\phi$ Accrual
algorithm. The implementation computes a dynamic suspicion value based on the elapsed time since the last successful
heartbeat and the historical distribution of heartbeat intervals.

\begin{lstlisting}[language=Java,
caption={Phi Accrual suspicion calculation},
label={lst:phi}]
double phi = phiDetector.calculatePhi(
    state.getSlidingWindowSeconds(),
    state.getLastAckTimeMs(),
    System.currentTimeMillis());

NodeStatus phiStatus = phiDetector.determineStatus(phi);
\end{lstlisting}

During the initial monitoring period, the detector may not have accumulated sufficient heartbeat observations to
perform reliable $\phi$ estimation. In the current implementation, the Phi Accrual Failure Detector requires at least
two historical heartbeat intervals before meaningful statistical analysis can be performed. To address this cold-start
problem, the system temporarily employs a conservative timeout-based strategy. Specifically, a node is not classified
as \textit{UNREACHABLE} unless it remains unresponsive for a duration substantially longer than the configured probe
interval. This grace period allows the detector to collect adequate heartbeat samples and establish a baseline of
normal communication behavior, thereby reducing the likelihood of false-positive failure detections for newly joined
nodes.

The computed $\phi$ value is mapped to one of four node states: \textit{ALIVE}, \textit{WARNING}, \textit{SUSPECTED},
or \textit{UNREACHABLE}. When a node transitions to either the \textit{SUSPECTED} or \textit{UNREACHABLE} state, the
information is disseminated through the gossip subsystem to ensure that neighboring nodes gradually converge toward a
consistent membership view.

\begin{lstlisting}[language=Java,
caption={Propagation of failure information},
label={lst:gossip}]
gossipService.gossipSuspect(targetNode);

gossipService.gossipUnreachable(targetNode);
\end{lstlisting}

If a node is ultimately classified as \textit{UNREACHABLE}, additional recovery procedures are triggered. The failed
node is removed from the local membership view, topology information is updated, and the failure event is reported to
the monitoring dashboard. Furthermore, the self-healing topology subsystem is notified to restore missing network
connections. Through the combination of periodic heartbeat monitoring, adaptive statistical failure estimation, gossip
dissemination, and automated recovery, SWIMPHI provides scalable and resilient failure detection for dynamic
distributed environments.

\subsubsection{FINDING NEMO}
\label{subsubsec:node_discovery_code}

This section highlights several key implementation aspects of FINDING NEMO. One important component is the construction
of the \textit{Join Plan}. Recall that acknowledgements received by a joining node can be classified into two
categories according to the state of the responder's neighbor list: \textit{fullNeighborAcks} and
\textit{missingNeighborAcks}. As discussed in Subsection~\ref{subsec:node_discovery}, a fully connected node must evict
one of its existing neighbors before accepting a new connection, whereas an under-connected node can directly
accommodate the joining node. Consequently, the implementation represents a \textit{JoinPlan} as a record containing
two components: a list of \textit{directTargets}, representing nodes that will establish direct connections with the
joining node, and a mapping \textit{evictionByDirectTarget}, which associates each fully connected target with its
selected eviction victim. The plan is generated by partitioning acknowledgements into the two categories and
determining the number of targets to be selected from each group.

\begin{lstlisting}[language=Java,
caption={Compute the quantity taken from two lists},
label={lst:quantity_target_list}]
int missingTargetCount;
int fullTargetCount;

int missingCount = missingNeighborAcks.size();
int fullCount = fullNeighborAcks.size();

if (missingCount + 2 * fullCount <= maxNeighbors) {
    missingTargetCount = missingCount;
    fullTargetCount = fullCount;
} else {
    int evenMissingCount = missingCount % 2 == 0 ? missingCount : missingCount - 1;
    evenMissingCount = Math.max(evenMissingCount, 0);

    missingTargetCount = Math.min(evenMissingCount, maxNeighbors);
    fullTargetCount = Math.min((maxNeighbors - missingTargetCount) / 2, fullCount);

    int minimumFullTargets = Math.min(Constant.MINIMUM_FULL_TARGET_COUNT, fullCount);

    if (fullTargetCount < minimumFullTargets && maxNeighbors >= 2 * minimumFullTargets) {
        fullTargetCount = minimumFullTargets;
        missingTargetCount = maxNeighbors - 2 * fullTargetCount;
        missingTargetCount = Math.max(0, missingTargetCount);
    }
}

List<JoinAck> directTargetAcks = new ArrayList<>();
directTargetAcks.addAll(missingNeighborAcks.subList(0, missingTargetCount));
directTargetAcks.addAll(fullNeighborAcks.subList(0, fullTargetCount));
List<JoinAck> chosenSufficientTargetAcks = sufficientNodeAcks.subList(0, sufficientTargetCount);
\end{lstlisting}

Listing~\ref{lst:quantity_target_list} illustrates the target-selection stage of FINDING NEMO. The algorithm
prioritizes nodes whose neighbor lists have not yet reached the target size, as these nodes can directly establish new
connections without requiring any topology reconfiguration. Such nodes are selected first and contribute one neighbor
each to the joining node. The remaining neighbor budget, if any, is then allocated to fully connected nodes. Since each
fully connected node contributes both itself and one eviction victim, selecting a single full-neighbor target
effectively provides two neighbors to the joining node. Consequently, the number of selected under-connected nodes is
rounded down to the nearest even value, ensuring that the remaining capacity can be satisfied through the eviction
mechanism while maintaining the desired target degree. Furthermore, the implementation attempts to include at least
\verb|MINIMUM_FULL_TARGET_COUNT| fully connected nodes whenever possible. By forcing new nodes to establish connections
with well-integrated regions of the overlay, the algorithm reduces the likelihood that bursts of simultaneous join
operations create isolated subnetworks and consequently mitigates the risk of network partitioning.

Another noteworthy implementation detail is the selection of eviction victims.
Listing~\ref{lst:evicted_members_selection} presents the corresponding procedure. For each fully connected direct
target, a neighboring node is selected uniformly at random as a potential eviction victim. An exclusion set is
maintained to prevent the joining node, previously selected direct targets, and previously selected victims from being
chosen again. This guarantees that each eviction victim is unique and cannot simultaneously serve as a direct target,
thereby avoiding conflicting topology updates.

\begin{lstlisting}[language=Java,
caption={Evicted members planning},
label={lst:evicted_members_selection}]
Set<String> excludedNodeIds = new HashSet<>();
excludedNodeIds.add(localAddress.nodeId());
excludedNodeIds.addAll(directTargetIds);

for (JoinAck directTarget : fullDirectTargetAcks) {
    List<NodeAddress> candidateVictims = new ArrayList<>(directTarget.responderNeighbors());
    Collections.shuffle(candidateVictims);
    for (NodeAddress victim : candidateVictims) {
        if (victim == null) continue;
        if (excludedNodeIds.contains(victim.nodeId())) continue;

        evictionByDirectTarget.put(directTarget.responder(), victim);
        excludedNodeIds.add(victim.nodeId());

        break;
    }
}
\end{lstlisting}

Finally, the protocol flow of FINDING NEMO is discussed. Due to the substantial amount of implementation-specific
boilerplate code, source-code listings are not used as the primary visualization aid. Instead, the protocol is
illustrated using the activity diagram shown in Figure~\ref{fig:FINDING_NEMO_protocol}, which more clearly conveys the
interaction sequence among participating nodes.

As shown in the diagram, the protocol diverges depending on whether the selected target node has available neighbor
capacity or a full neighbor list. Despite this distinction, the two execution paths share a largely similar sequence of
operations, differing primarily in the number of neighbors ultimately acquired by the joining node. The diagram also
highlights that a target node communicates with the selected eviction victim and completes the corresponding topology
update before transmitting a \texttt{COMMIT\_ACK} message to the joining node.

An earlier version of the protocol employed a more aggressive update strategy. Upon receiving a \texttt{COMMIT\_DIRECT}
message, the target node simultaneously added the joining node and removed the selected victim. Subsequently, a
\texttt{COMMIT\_VICTIM} message instructed the eviction victim to add the joining node and remove the target node.
Although conceptually straightforward, this approach introduced the possibility of partial commits. In particular,
message loss or node failure during the update sequence could leave one endpoint retaining a neighbor relationship
while the corresponding peer had already removed it, resulting in an inconsistent overlay topology.

To mitigate this issue, the current protocol adopts a more conservative commit procedure, ensuring that topology
updates are coordinated before acknowledgements are returned. While temporary degree imbalance may still occur, such
situations are subsequently corrected by the rewiring mechanism described later in this work.

\begin{figure*}[!t]
  \centering
  \includegraphics[width=0.95\textwidth]{FINDING_NEMO_actiDi.png}
  \caption{FINDING NEMO protocol flow}
  \label{fig:FINDING_NEMO_protocol}
\end{figure*}

\subsubsection{CHUALANH}
\label{subsubsec:node_deletion_code}
\begin{enumerate}
  \item{\textbf{Direct Rewiring: }}

        The code fragment in Listing~\ref{lst:direct_rewiring} presents the core implementation of the direct rewiring
        procedure for \textbf{Case 1: No Existing Connection}, described in Section~\ref{subsec:node_deletion}. Each direct
        rewiring session is uniquely identified by a recovery identifier generated using a \texttt{UUID}. The initiating
        deficient node, denoted as node(A), constructs a \texttt{REWIRE\_REQ\_DIRECT} message and sends it to another deficient
        node, denoted as node(B).

        Upon receiving the request, node (B) evaluates whether it can participate in the rewiring process. If the request is
        accepted, node (B) returns a response message with the status \texttt{ACCEPTED}. After receiving the acceptance, node
        (A) performs the actual connection establishment by invoking the \texttt{connectionManager.addIfSpace()} method, which
        adds node (B) to its neighbor list provided that sufficient capacity is available.

        Subsequently, node (A) sends a \texttt{REWIRE\_DIRECT\_COMMIT} message to node (B) to finalize the rewiring session.
        Upon processing this commit message, node (B) similarly invokes \texttt{connectionManager.addIfSpace()} to add node (A)
        to its neighbor list. Consequently, a bidirectional connection is established between the two deficient nodes,
        restoring one missing neighbor for each participant while maintaining the degree constraints imposed by the network
        topology.

        \begin{lstlisting}[language=Java,
    caption={Direct Rewiring},
    label={lst:direct_rewiring}]
private boolean tryDirectRepair(
    NodeAddress defA, NodeAddress defB) {
    String recoveryId =
        UUID.randomUUID().toString();

    RewireMessage request = RewireMessage.of(
        RecoveryMessageType.REWIRE_REQ_DIRECT,
        recoveryId, localAddress,
        defA, defB, null, null, null, null,
        connectionManager.neighborAddresses(),
        List.of(), null);

    RewireMessage response = null;
    try {
        response = sendAsync(defB, request)
            .get(REQUEST_TIMEOUT.toMillis(),
                TimeUnit.MILLISECONDS);
    } catch (InterruptedException
        | ExecutionException
        | TimeoutException e) {
        Console.logError("[REWIRE] An unexpected error has occured " + e.getMessage());    
        return false;
    }

    if (response == null || response.status()
            != RewireStatus.ACCEPTED) {
        Console.log("[REWIRE] response is null or not ACCEPTED");
        return false;
    }

    boolean localApplied = connectionManager
        .addIfSpace(defB,
            "direct rewiring recovery commit");

    if (!localApplied) {
        Console.log("[REWIRE] Local applied failed");
        return false;
    }

    RewireMessage commit = RewireMessage.of(
        RecoveryMessageType.REWIRE_DIRECT_COMMIT,
        recoveryId, localAddress,
        defA, defB, null, null, defB, null,
        connectionManager.neighborAddresses(),
        List.of(), RewireStatus.ACCEPTED);

    sendOnly(defB, commit);
    refreshHealthState();
    Console.logSuccess("Direct Repait completed successfully");
    return true;
}
\end{lstlisting}

  \item{\textbf{Undirect Rewiring: }}

        The code in \texttt{runFullRewiring()} method presents the core implementation of the rewiring procedure for
        \textbf{Case 2: Existing Connection}, in which two deficient nodes are already directly connected. Initially, the
        current deficient node, denoted as node (A), discovers another deficient node, denoted as node (B), with which it
        already maintains a direct connection. Node (A) initiates a rewiring session by generating a unique request identifier
        and transmitting a \texttt{REWIRE\_SESSION\_REQ} message to node (B).

        \begin{lstlisting}[language=Java,
    caption={Generating Session ID},
    label={lst:sessionId}]
String requestId = UUID.randomUUID().toString();

RewireMessage req = RewireMessage.of(
    REWIRE_SESSION_REQ,
    requestId, localAddress,
    defA, defB, ...);

RewireMessage ack = sendAsync(defB, req)
    .get(REQUEST_TIMEOUT.toMillis(),
    TimeUnit.MILLISECONDS);

    \end{lstlisting}

        Upon receiving the acknowledgment from node (B), a leader election mechanism is executed to determine which deficient
        node will coordinate the recovery process. Without loss of generality, assume that node (A) is elected as the leader.
        Node (A) subsequently generates a recovery identifier and sends a \texttt{REWIRE\_SESSION\_COMMIT} message to node (B),
        informing node (B) that it will participate in the recovery process as a fellow.

        \begin{lstlisting}[language=Java,
    label={lst:fellow}]
if (!localAddress.nodeId().equals(
    electLeader(defA, defB).nodeId())) {
    return false;
}

String recoveryId = UUID.randomUUID().toString();

sendOnly(defB, RewireMessage.of(
    REWIRE_SESSION_COMMIT,
    recoveryId,
    localAddress,
    defA, defB, ...));

    \end{lstlisting}

        The leader then computes a set of candidate nodes, denoted by (C), where each node (C) is a neighbor of node (A) but is
        not directly connected to node (B). This candidate set is obtained by examining the entire neighbor list of node (A)
        and filtering out all nodes that also appear in the neighbor list of node (B). For each candidate (C), the leader sends
        a query message containing neighbor lists of node A and B, along with the recovery status. Upon receiving the query,
        node (C) searches for a node (D) that is a neighbor of (C) but not a neighbor of node (A). Once such a pair ((C,D)) is
        identified, the CHUALANH algorithm is executed to restructure the topology while preserving node degrees.
        \begin{lstlisting}[language=Java,
    caption={Finding Node C and D},
    label={lst:findcd}]
List<NodeAddress> candidateCs = selectCandidateCs( defANeighbors, defBNeighbors);
for (NodeAddress c : candidateCs) {
    RewireMessage response = sendAsync(c, query).get(...);
    if (response != null && response.nodeD() != null) {
        NodeAddress d = response.nodeD();
        if (executeScheme( recoveryId, defA, defB, c, d)) {
            return true;
        }
    }
}
    \end{lstlisting}

        If such a node (D) exists, node (A) invokes the \texttt{executeScheme()} method, which applies the \textbf{Connectivity
          Healing through Under-connected and Adaptive Link Adjustment for Network Health (CHUALANH)} for \textbf{Case 2:
          Existing Connection}, as described in Section~\ref{subsec:node_deletion}. The algorithm restructures the network by
        replacing selected edges while preserving the degree constraints of the participating nodes. Consequently, the
        deficient nodes recover their missing connections and the topology is repaired without altering the predefined node
        degree (k)
\end{enumerate}

\subsection{Server Application}
\label{subsec:server_application}
The monitoring dashboard follows a client-server architecture in which each distributed node in decentralized
network acts as a client and the dashboard application serves as a centralized server reponsible for
collecting, storing, and visualizing network information

\subsubsection{\textbf{Node Status Reporting}}
\label{subsubsec: nodestatusreport}
Each node contains a \mbox{\texttt{DashboardReporter}} component that periodically reports its current state
to the dashboard. Communication between nodes and the dashboard is implemented using HTTP and JSON messages.

When a node is operational, the \mbox{\texttt{reportSelfAlive()}} method constructs a JSON object containing the node's
id (\mbox{\texttt{id}}), IP address(\texttt{ipAddress}), peer-to-peer communication port(\texttt{p2pPort}),
password-cracking service port(\texttt{crackingPort}), current neighbor list (\texttt{neighbors}), and a status field
set to \texttt{"UP"}. The message is then transmitted asynchronously to the dashboard's \texttt{/hearbeat} endpoint
using and HTTP POST request.
\begin{lstlisting}[language=Java,
    label={lst:evicted_members_selection}]
String json = """
    {
        "id": "%s",
        "ipAddress": "%s",
        "p2pPort": %d,
        "crackingPort": %d,
        "status": "UP",
        "neighbors": [%s]
    }
    """.formatted(localNodeId, advertiseHost, p2pPort, crackingPort, neighborsJson);

HttpRequest request = HttpRequest.newBuilder()
    .uri(URI.create(dashboardUrl + "/heartbeat"))
    .timeout(Duration.ofSeconds(5))
    .header("Content-Type", "application/json")
    .POST(HttpRequest.BodyPublishers.ofString(json))
    .build();

    \end{lstlisting}

In the event of a node failure, a neighboring node that detects the communication loss reports the failure through the
\texttt{reportFailure()} method. This method generates a JSON message containing the reporting node identifier
(\texttt{reporterNodeId}), the suspected failed node identifier (\texttt{failedNodeId}), the computed phi value
(\texttt{phi}), the failure detection threshold (\texttt{threshold}), a timestamp (\texttt{timestamp}), and a status
field set to "\texttt{UNREACHABLE}". The message is subsequently sent to the dashboard's \texttt{/failure-report}
endpoint using an HTTP POST request. The use of asynchronous HTTP requests through \texttt{CompletableFuture} ensures
that dashboard communication does not block the execution of the distributed system.
\begin{lstlisting}[language=Java,
    label={lst:evicted_members_selection}]
String json = String.format(Locale.US, """
        {"reporterNodeId": "%s",
        "failedNodeId": "%s",
        "message": "%s",
        "phi": %.6f,
        "threshold": %.6f,
        "status": "UNREACHABLE",
        "timestamp": "%s"
        }""",
    localNodeId,
    failedNode.nodeId(),
    message,
    phi,
    threshold,
    LocalDateTime.now());

HttpRequest request = HttpRequest.newBuilder().uri(URI.create(dashboardUrl + "/failure-report")).timeout(Duration.ofSeconds(5))
.header("Content-Type", "application/json")
.POST(HttpRequest.BodyPublishers.ofString(json)).build();
\end{lstlisting}
\subsubsection{\textbf{Server-Side Request Processing}}
\label{subsubsec: server}
On the server side, the \texttt{NodeController} class exposes REST endpoints that receive and process
requests from distributed nodes. Incoming hearbeat messages are mapped to the \texttt{/heartbeat} endpoint
and handled by the \texttt{receiveHeartbeat()} method. This method delegates the processing task to
\texttt{NodeService.processHeartbeat()}, which updates or creates anode record in the database. The service
stores the node's network information, neighbor list, and latest heartbeat timestamp while setting its status
to "\texttt{UP}". Similarly, failure notifications are mapped to the \texttt{/failure-report} endpoint and
handled by the \texttt{receiveFailureReport()} method. The method invokes
\texttt{NodeService.processFailureReport()} to validate and persist the failure report. Subsequently, the
reported is marked as "\texttt{FAILED}" in the node repository, and the failure information is stored in the
failure report repository for later visualization and analysis.

\begin{lstlisting}[language=Java,
label={lst:evicted_members_selection}]
@PostMapping("/heartbeat")
public ResponseEntity<String> receiveHeartbeat(
        @RequestBody Node node) {
    nodeService.processHeartbeat(node);
    return ResponseEntity.ok(
            "Heartbeat received from " + node.getId());
}

@PostMapping("/failure-report")
public ResponseEntity<String> receiveFailureReport(
        @RequestBody FailureReport report) {
    nodeService.processFailureReport(report);
    passwordCrackingService.handleNodeFailure(report.getFailedNodeId());
    return ResponseEntity.ok(
            "Failure report received");
}
\end{lstlisting}

\subsubsection{\textbf{Dashboard Visualization}}
\label{subsubsec: dashboard}
The dashboard obtains monitoring information through REST endpoints exposed by the server. The
\texttt{/nodes} and \texttt{/nodes/{id}} endpoints return information about all nodes or a specific node,
respectively, while the \texttt{/failure-reports} endpoint returns the history of detected failures. The
dashboard periodically issues HTTP GET requests to these endpoints to retrieve the latest data stored in the
repositories and visualize the current system state, including node availability, network topology, and
recorded failures.
\begin{lstlisting}[language=Java,
label={lst:evicted_members_selection}]
@GetMapping("/nodes")
public ResponseEntity<List<Node>> getAllNodes() {
    return ResponseEntity.ok(nodeService.getAllNodes());
}

@GetMapping("/nodes/{id}")
public ResponseEntity<Node> getNodeById(@PathVariable String id) {
    Node node = nodeService.getNodeById(id);

    if (node == null) {
        return ResponseEntity.notFound().build();
    }

    return ResponseEntity.ok(node);
}

@GetMapping("/failure-reports")
public ResponseEntity<List<FailureReport>> getFailureReports() {
    return ResponseEntity.ok(nodeService.getFailureReports());
}
\end{lstlisting}
